\subsubsection{Layer 3 - Core Domain (Entities \& Value Objects)}

\paragraph{Cliente}
\textbf{Descrizione della classe:} rappresenta il cliente commerciale autenticato e abilitato all'inserimento ordini.

\textbf{Descrizione degli attributi:}
\begin{itemize}
    \item \texttt{id}: identificativo univoco del cliente.
    \item \texttt{ragioneSociale}: denominazione legale dell'azienda cliente.
    \item \texttt{email}: contatto email principale per autenticazione/comunicazioni.
\end{itemize}

\textbf{Descrizione dei metodi:} non sono definiti metodi nel diagramma per questa entita.

\begin{figure}[H]
    \centering
    \includegraphics[width=0.48\textwidth]{uml_backend/14_cliente.png}
    \caption{UML della classe \texttt{Cliente}.}
\end{figure}

\paragraph{Operatore}
\textbf{Descrizione della classe:} identifica l'utente interno che prende in carico ticket e supervisiona i casi critici.

\textbf{Descrizione degli attributi:}
\begin{itemize}
    \item \texttt{id}: identificativo univoco operatore.
    \item \texttt{nome}: nominativo operatore.
\end{itemize}

\textbf{Descrizione dei metodi:} non sono definiti metodi nel diagramma.

\begin{figure}[H]
    \centering
    \includegraphics[width=0.45\textwidth]{uml_backend/15_operatore.png}
    \caption{UML della classe \texttt{Operatore}.}
\end{figure}

\paragraph{Order}
\textbf{Descrizione della classe:} rappresenta l'ordine consolidato pronto per storicizzazione ed esportazione verso ERP.

\textbf{Descrizione degli attributi:}
\begin{itemize}
    \item \texttt{id}: identificativo ordine.
    \item \texttt{clienteId}: riferimento al cliente proprietario dell'ordine.
    \item \texttt{dataOrdine}: timestamp di emissione.
    \item \texttt{righe}: collezione delle righe ordine derivate dal carrello.
\end{itemize}

\textbf{Descrizione dei metodi:} nel diagramma non sono mostrati metodi per questa entita.

\begin{figure}[H]
    \centering
    \includegraphics[width=0.52\textwidth]{uml_backend/16_order.png}
    \caption{UML della classe \texttt{Order}.}
\end{figure}

\paragraph{Carrello}
\textbf{Descrizione della classe:} modella lo stato temporaneo dell'ordine in costruzione per una sessione chat.

\textbf{Descrizione degli attributi:}
\begin{itemize}
    \item \texttt{chatsessionId}: identificativo sessione conversazionale associata.
    \item \texttt{clienteId}: cliente a cui il carrello appartiene.
    \item \texttt{righe}: elenco righe prodotto correnti.
\end{itemize}

\textbf{Descrizione dei metodi:}
\begin{itemize}
    \item \texttt{calcolaTotalePezzi()}: restituisce la somma totale delle quantita nel carrello.
\end{itemize}

\begin{figure}[H]
    \centering
    \includegraphics[width=0.54\textwidth]{uml_backend/17_carrello.png}
    \caption{UML della classe \texttt{Carrello}.}
\end{figure}

\paragraph{RigaCarrello}
\textbf{Descrizione della classe:} rappresenta una singola posizione prodotto-quantita nel carrello.

\textbf{Descrizione degli attributi:}
\begin{itemize}
    \item \texttt{prodottoId}: riferimento all'articolo di catalogo.
    \item \texttt{quantita}: numero unita richieste.
    \item \texttt{confidenceScore}: affidabilita della proposta AI su quella riga.
\end{itemize}

\textbf{Descrizione dei metodi:} non sono definiti metodi nel diagramma.

\begin{figure}[H]
    \centering
    \includegraphics[width=0.50\textwidth]{uml_backend/18_rigacarrello.png}
    \caption{UML della classe \texttt{RigaCarrello}.}
\end{figure}

\paragraph{Session}
\textbf{Descrizione della classe:} mantiene il contesto di conversazione tra cliente e sistema durante la composizione ordine.

\textbf{Descrizione degli attributi:}
\begin{itemize}
    \item \texttt{id}: identificativo sessione.
    \item \texttt{clientId}: identificativo del cliente associato.
    \item \texttt{status}: stato della sessione (aperta, chiusa, sospesa).
    \item \texttt{messages}: storico messaggi scambiati.
\end{itemize}

\textbf{Descrizione dei metodi:} non sono presenti metodi nel diagramma.

\begin{figure}[H]
    \centering
    \includegraphics[width=0.50\textwidth]{uml_backend/19_session.png}
    \caption{UML della classe \texttt{Session}.}
\end{figure}

\paragraph{Ticket}
\textbf{Descrizione della classe:} rappresenta una richiesta di supporto aperta per intervento umano.

\textbf{Descrizione degli attributi:}
\begin{itemize}
    \item \texttt{id}: identificativo ticket.
    \item \texttt{clienteId}: cliente che ha aperto il ticket.
    \item \texttt{stato}: stato corrente del ticket.
    \item \texttt{dataApertura}: data e ora di apertura.
\end{itemize}

\textbf{Descrizione dei metodi:} il diagramma non esplicita metodi per questa entita.

\begin{figure}[H]
    \centering
    \includegraphics[width=0.48\textwidth]{uml_backend/20_ticket.png}
    \caption{UML della classe \texttt{Ticket}.}
\end{figure}

\paragraph{LogAI}
\textbf{Descrizione della classe:} traccia input, output e correzioni legate alle elaborazioni AI per audit e miglioramento modello.

\textbf{Descrizione degli attributi:}
\begin{itemize}
    \item \texttt{id}: identificativo del log.
    \item \texttt{inputOriginario}: testo/audio trascritto ricevuto dal cliente.
    \item \texttt{rispostaLLM}: proposta generata dal modello linguistico.
    \item \texttt{correzioneDelta}: differenza applicata manualmente da utente/operatore.
    \item \texttt{stato}: stato del log nel workflow di validazione.
\end{itemize}

\textbf{Descrizione dei metodi:} non sono definiti metodi nel diagramma.

\begin{figure}[H]
    \centering
    \includegraphics[width=0.55\textwidth]{uml_backend/21_logai.png}
    \caption{UML della classe \texttt{LogAI}.}
\end{figure}
